Der Auftrieb ist die Gewichtskraft der verdrängten Luft, $\rho_\text{L}\,V\,g$. Er muss den leeren Zeppelin (Masse $m_\text{Z}$), das Wasserstoffgas (Masse $\rho_\text{H}\,V$) und die Nutzlast $m_\text{N}$ tragen:
\begin{align}
 \rho_\text{L}\,V\,g &= \left(m_\text{Z} + \rho_\text{H}\,V + m_\text{N}\right)g
\end{align}
Mit den Dichten bei Normalbedingungen, $\rho_\text{L} = \rhoLO$ und $\rho_\text{H} = \rhoHO$ (FoTaBe S.~187), wird die Nutzlast
\begin{align}
 m_\text{N} &= \mNF \\
   &= \left(\rhoL-\rhoH\right)\cdot\VZ-\mZ \\
   &= \mN \approx \result{\mNtP{0}{2}}
\end{align}
